\documentclass[a4paper]{article}

\usepackage{graphicx}
\usepackage{amsmath,amssymb}
\usepackage{minted}
\usepackage{multirow}
\usepackage{float}
\usepackage{todonotes}
\usepackage{forest}
\usepackage{xstring}
\usepackage{xcolor}
\usepackage[most]{tcolorbox}
\usepackage{xparse}
\usepackage{natbib}

\usetikzlibrary{graphs,graphdrawing}
\usegdlibrary{layered}

% for external graphviz
\input{/home/donkarlo/Dropbox/repo/nd_ai_project/src/nd_ai/knoweldge_representation/language/formal/computer/programming/tex/latex/code_clips/external_graphviz.tex}

% for tags
\input{/home/donkarlo/Dropbox/repo/nd_ai_project/src/nd_ai/knoweldge_representation/language/formal/computer/programming/tex/latex/parts/control_sequence/kind/semtag/semtag.tex}

% for links
\input{/home/donkarlo/Dropbox/repo/nd_ai_project/src/nd_ai/knoweldge_representation/language/formal/computer/programming/tex/latex/code_clips/seehere.tex}

\title{}
\date{}

\begin{document}
    \maketitle
    
    
    \section{zwar \dots aber}
        
        \subsection{Positiv}
            Das Handy ist zwar billig, aber gut: The phone is indeed cheap, but good. \\
            Ich habe zwar Zeit, aber Energie: I do have time, but (also) energy. \\
            Er ist zwar müde, aber motiviert: He is indeed tired, but motivated. \\
        
        \subsection{Negativ}
            Das Handy ist zwar billig, aber nicht gut: The phone is cheap, but not good. \\
            Ich habe zwar Zeit, aber keine Energie: I do have time, but no energy. \\
            Er ist zwar müde, aber nicht motiviert: He is tired, but not motivated. \\
    
    
    \section{allerdings}
        
        \subsection{Positiv}
            Das ist eine gute Idee, allerdings kompliziert: That is a good idea, however complicated. \\
            Ich komme gern, allerdings später: I would like to come, however later. \\
            Er ist freundlich, allerdings direkt: He is friendly, however direct. \\
        
        \subsection{Negativ}
            Das ist eine gute Idee, allerdings nicht einfach: That is a good idea, however not easy. \\
            Ich komme gern, allerdings nicht heute: I would like to come, however not today. \\
            Er ist freundlich, allerdings nicht immer: He is friendly, however not always. \\
    
    
    \section{freilich}
        
        \subsection{Positiv}
            Das ist freilich möglich: That is of course possible. \\
            Er hat freilich recht: He of course is right. \\
            Wir können freilich gehen: We can of course go. \\
        
        \subsection{Negativ}
            Das ist freilich nicht einfach: That is of course not easy. \\
            Er hat freilich nicht recht: He of course is not right. \\
            Wir können freilich nicht bleiben: We of course cannot stay. \\
    
    
    \section{gewiss}
        
        \subsection{Positiv}
            Das ist gewiss richtig: That is certainly correct. \\
            Er kommt gewiss morgen: He will certainly come tomorrow. \\
            Sie ist gewiss bereit: She is certainly ready. \\
        
        \subsection{Negativ}
            Das ist gewiss nicht richtig: That is certainly not correct. \\
            Er kommt gewiss nicht morgen: He will certainly not come tomorrow. \\
            Sie ist gewiss nicht bereit: She is certainly not ready. \\
    
    
    \section{schon}
        
        \subsection{Positiv}
            Das wird schon gut: It will be fine. \\
            Er hat schon recht: He is indeed right. \\
            Das passt schon: That is fine. \\
        
        \subsection{Negativ}
            Das wird schon nicht schlecht: It will not be bad. \\
            Er hat schon nicht unrecht: He is not entirely wrong. \\
            Das passt schon nicht: That does not really fit. \\
    %\bibliographystyle{plainnat}
    %\bibliography{/home/donkarlo/Dropbox/repo/nd_shared_research/bibliography/bibliography.bib}
\end{document}